Fix an integer n>=4. Let F_n be the Farey sequence of order n: the list of all reduced fractions a/b in [0,1] with 0<=a<=b<=n and gcd(a,b)=1, arranged in increasing numerical order. Write this sequence as
  a_1/b_1 < a_2/b_2 < ... < a_M/b_M.
Define f(n) to be the largest integer f with the property:
  for all indices 1<=k<l<=M with l-k <= f, one has (a_k - a_l)(b_k - b_l) >= 0.
Equivalently: any two Farey fractions within index-distance <= f are "similarly ordered", meaning the ordering of numerators agrees with the ordering of denominators (allowing ties).
Question: estimate f(n). In particular, does there exist a constant c>0 such that f(n) = (c+o(1)) n?
